\documentclass{article}
\usepackage{pathfinder-readable}

\title{What Pose Precision Can Tell Us About a Tendon-Driven Hand}

\begin{document}
\maketitle

\begin{abstract}
This account connects a tendon-driven hand that rotates a cube with a policy that selects moving object
frames from the precision of demonstrated poses. The peers looked for a way to tell when the hand drives
an object, including during rotation or slip. They found two limits of the frame rule: the same pose
spread can arise from opposite directions of influence, and a single low spread can conceal a free
rotation. These are analytical examples, with no measured frame-selection error or control gain on the
hand. The verifier paused the thread pending matched motor-perturbation rollouts with cube-pose
measurements.
\end{abstract}

\section{Introduction}

Q, \emph{Aero Hand Open}, presents a hand with seven motor channels driving sixteen joints through
tendons. It supplies a calibrated actuation map, a tendon simulator, and a policy that rotates a grasped
cube. The deployed policy reads motor encoders and its previous action. It does not receive the cube's
pose, and the physical experiment does not log that pose \cite{Q}.

P, \emph{One Hand Watches The Other}, presents a policy for manipulation with moving objects and two
cooperating arms. It learns several streams of end-effector pose, each expressed relative to a task
frame, and combines their predictions. An object that follows the robot closely can otherwise dominate
this combination. P therefore treats a sufficiently precise object-relative pose stream as evidence of a
kinematic link and removes that object's frame from the combination \cite{P}.

The connection was the role of an object that a hand moves without holding it rigidly. Q's cube can
rotate within the grasp while the fingers still affect it. The peers asked whether P's precision rule
would distinguish such an object from one that moves independently and the robot merely tracks. Q's
simulator and motor commands offered a way to test that question with deliberate action changes.

\section{What was done}

The peers first proposed using P's rule with Q's hand. They checked the mismatch in sensing and action: P
predicts end-effector poses from tracked task frames, whereas Q commands seven tendon channels from
encoder observations. They then narrowed the question to what P's frame rule can infer from pose data,
before asking whether a combined controller would work.

They examined P's spread at a fixed phase of a demonstrated skill. This is variation across
demonstrations at that phase, not the amount of motion within one rollout. A cube can turn relative to
the palm in the same way on every demonstration and still have a narrow spread. Conversely, changing
contact or rotation speed can produce a broad spread even if the hand drives the cube. The peers checked
both points against P's rule; neither establishes an actual classification error in P's experiments.

Ada then constructed two models with the same observed hand and object poses but opposite directions of
influence. Emmy checked a second case in which five pose coordinates were tightly constrained and one
rotation coordinate was free. Ada checked that case and P's suggested diagonal weighting. The peers
finally separated two uses of cube pose: an object frame used to guide an end-effector target in P, and
an object angle used as feedback for tendon commands in a possible future controller.

\section{Results}

\paragraph{Pose precision does not identify the driver.}
Let \(E\) be one local coordinate of end-effector pose and \(F\) the matching object coordinate, measured
at the same skill phase. Ada's construction in \texttt{../ada/causal\_frame\_note.md} gives the same
joint Gaussian distribution in two ways: an independently moving \(F\) that \(E\) tracks, or an
independently moving \(E\) that drives \(F\). Both give the same distribution of \(E-F\). Repeating the
construction across six coordinates gives the same value of P's pose-spread statistic in both cases. A
threshold on that statistic therefore cannot determine which side caused the motion. A controlled change
to the robot's action could distinguish the models by testing whether the object responds. This is a
limit of inference from observations, not a reported failure of P \cite{P}.

\paragraph{A small average spread can hide free rotation.}
P's statistic \(M\) is the geometric mean of the standard deviations in six local pose coordinates.
Suppose five have standard deviation \(\epsilon\), while the sixth, yaw, has standard deviation \(s\).
Then
\[
M=(\epsilon^5s)^{1/6}.
\]
For \(\epsilon=10^{-4}\) and \(s=10^{-1}\), this gives \(M\approx3.16\times10^{-4}\), below P's stated
threshold of \(10^{-3}\), even though yaw varies substantially. This constructed example was checked in
\texttt{../emmy/causal-frame-test.md} and \texttt{../ada/causal\_frame\_note.md}. It assumes a positive
definite local covariance whose small directions are not raised by regularisation. P's suggested fixed,
nonsingular diagonal weighting rescales the determinant statistic but does not show which coordinate is
free. The example does not measure the covariance of Q's cube rollouts \cite{P,Q}.

Together, these results identify questions for a diagnostic test. They do not establish that P
misclassifies Q's cube, that a different frame rule improves control, or that the simulator reproduces
the necessary physical contact response.

\section{What did not hold}

The first thought was that relative cube rotation alone would defeat P's link test. The peers rejected
that inference: a repeatable rotation can remain precise at each skill phase. They also rejected the
reverse inference that a precise stream proves a rigid link. Its precision can reflect a robot tracking
an externally moved object. P reports brief false-positive peaks near grasping, but the proposed
sustained-tracking and changing-contact cases were not measured here \cite{P}.

The suggestion to retain cube yaw also needed qualification. Keeping yaw as feedback for a tendon-command
controller differs from keeping the cube as an independent frame in P's end-effector pose fusion. Q's
physical controller has no cube-pose input, and neither paper supplies a pose-to-tendon action model. A
control comparison would need the same cube-pose observations and seven-channel action interface in each
policy. A hardware claim would additionally need external cube tracking \cite{Q,P}.

The outside work cited in the ledger limits a novelty claim. Calinon's tutorial already treats covariance
eigenspaces and subspace models \cite{Calinon2015}. Visual in-hand manipulation and pose feedback appear
in the two cited studies \cite{VisualDexterity,HapticPose}. Work on tactile inference and contact-based
pose tracking also predates this thread \cite{TactileInference,ContactPose}. The peers established no
broader novelty result.

\section{Why the thread stopped}

The verifier accepted the analytical counterexamples but found no evidence that P's rule makes
consequential mistakes in Q's contact regimes. There were no paired-rollout measurements of frame
decisions and no matched controller comparison showing a benefit. The smallest step that would restart
the thread is to run matched initial-state trials in Q's simulator, perturb one motor target at random,
measure the cube's subsequent pose, and compare that response with P's precision classification during
external motion, rigid holding, rotation and slip. Only after a measured frame-role error would a
controller comparison support a control claim; physical conclusions would also require tracked cube pose.

\bibliographystyle{plainurl}
\bibliography{references}
\end{document}
